\documentclass{beamer}

\usetheme{Berlin}
\usecolortheme{Orchid}
\useoutertheme{miniframes}
\setbeamertemplate{navigation symbols}{}

\usepackage[utf8]{inputenc}
\usepackage[T1]{fontenc}
\usepackage{amsmath}
\usepackage{amssymb}
\usepackage{booktabs}
\usepackage{graphicx}
\graphicspath{{../images/}}

\usepackage{xcolor}
\usepackage{listings}
\usepackage{hyperref}
\usepackage{tikz}
\usetikzlibrary{arrows.meta,positioning,shapes.geometric,calc}

\lstset{
  language=Java,
  basicstyle=\ttfamily\small,
  keywordstyle=\color{blue},
  commentstyle=\color{gray},
  stringstyle=\color{teal},
  showstringspaces=false,
  columns=fullflexible,
  breaklines=true
}

% Per-slide source line (references shown on the slide itself).
\newcommand{\slidesources}[1]{%
  \vfill
  {\tiny\color{gray}\raggedright\sloppy\parbox{\linewidth}{Sources: #1}\par}%
}

% Unit-wise source strings for this subject (used by slides/notes).
% Path is relative to the lecture `latex/` folder (the compilation CWD).
% OOPS with Java (CSEG1044) — unit-wise sources
%
% These are short, student-facing reference strings intended to be used with:
%   - \slidesources{...}  (in beamer slides)
%   - \notesources{...}   (in lecture notes)
%
% Style inspiration: other subjects in My-Subjects/ use semicolon-separated sources
% and include an "accessed" date for web documentation.

\newcommand{\OOPSUnitISources}{Schildt, \textit{Java: A Beginner's Guide} (10e), McGraw-Hill (Java fundamentals + OOP basics).; Oracle Java Tutorials: Getting Started + Language Basics + Classes and Objects (accessed \today).; Java SE API docs (e.g., \texttt{String}, \texttt{Scanner}) (accessed \today).}

\newcommand{\OOPSUnitIISources}{Schildt, \textit{Java: The Complete Reference} (12e), McGraw-Hill (classes, inheritance, packages, interfaces).; Bloch, \textit{Effective Java} (3e), Addison-Wesley, 2018 (design choices; interfaces vs abstract classes).; Oracle Java Tutorials: Classes and Objects + Interfaces + Packages (accessed \today).; Java SE API docs (e.g., \texttt{Comparable}, \texttt{Runnable}) (accessed \today).}

\newcommand{\OOPSUnitIIISources}{Schildt, \textit{Java: The Complete Reference} (12e) (nested/inner classes, exceptions, threads, I/O).; Oracle Java Tutorials: Nested Classes + Exceptions + Concurrency + Basic I/O (accessed \today).; Java SE API docs (e.g., \texttt{Thread}, \texttt{AutoCloseable}) (accessed \today).}

\newcommand{\OOPSUnitIVSources}{Schildt, \textit{Java: The Complete Reference} (12e) (generics, lambdas, Swing, JDBC).; Bloch, \textit{Effective Java} (3e) (generics + lambdas).; Oracle Java Tutorials: Generics + Lambda Expressions + Swing + JDBC Basics (accessed \today).; Java SE API docs (e.g., \texttt{java.util.function}, \texttt{javax.swing}, \texttt{java.sql}) (accessed \today).}

\newcommand{\OOPSUnitVSources}{Schildt, \textit{Java: The Complete Reference} (12e) (Collections Framework + wrapper classes).; Oracle Java docs: Collections Framework overview (accessed \today).; Java SE API docs (\texttt{java.util} collections; wrapper classes in \texttt{java.lang}) (accessed \today).}

\newcommand{\OOPSUnitVISources}{Git SCM: \textit{Pro Git} (2e) (version control workflow), accessed \today.; GitHub Docs (repositories, issues, pull requests), accessed \today.; Oracle Java Tutorials: Swing + JDBC (accessed \today).; JUnit 5 User Guide (accessed \today).; Apache Maven Project: Getting Started (accessed \today).; Gradle User Manual: Getting Started (accessed \today).; UML class diagrams: UML 2.x overview/spec (accessed \today).}

\newcommand{\OOPSMiscWhyInterfacesSources}{Schildt, \textit{Java: The Complete Reference} (12e) (interfaces).; Bloch, \textit{Effective Java} (3e) (prefer interfaces where appropriate).; Oracle Java Tutorials: Interfaces (accessed \today).; Java SE API docs (e.g., \texttt{Comparable}, \texttt{Runnable}) (accessed \today).}



\title[OOPS with Java]{Object Oriented Programming (CSEG1044)}
\subtitle{Unit 02 -- Lecture 02: Constructors, `this`, Overloading, `static`, and GC Overview}
\begin{document}

\begin{frame}[fragile]
  \titlepage
  \vspace{0.5em}
  \slidesources{\OOPSUnitIISources}
\end{frame}

\begin{frame}{Quick Links}
  \centering
  \hyperlink{sec:overview}{\beamerbutton{Overview}}\hspace{0.6em}
  \hyperlink{sec:concepts}{\beamerbutton{Key Topics}}\hspace{0.6em}
  \hyperlink{sec:demo}{\beamerbutton{Demo}}\hspace{0.6em}
  \hyperlink{sec:summary}{\beamerbutton{Summary}}
  \slidesources{\OOPSUnitIISources}
\end{frame}

\begin{frame}{Agenda}
  \tableofcontents
  \slidesources{\OOPSUnitIISources}
\end{frame}

\section{Overview}
\label{sec:overview}

\begin{frame}{Lecture Snapshot}
\begin{itemize}[<+->]
  \item \textbf{Unit focus:} Classes, Inheritance, Packages, and Interfaces
  \item \textbf{Course thread:} Use Constructors, `this`, Overloading, `static`, and GC Overview to strengthen the domain model, APIs, and access rules inside Campus Resource Manager.
\end{itemize}
  \slidesources{\OOPSUnitIISources}
\end{frame}

\begin{frame}{Recall Before You Start}
\begin{itemize}[<+->]
  \item \textbf{Class state and behavior}: Constructors and overloading only make sense once fields and methods already mean something in one class. Main reference: \texttt{Unit 02 / Lecture 01 slides/notes}.
  \item \textbf{Reference semantics}: Constructors make object creation explicit, not just primitive assignment. Main reference: \texttt{Unit 02 / Lecture 01 slides/notes}.
\end{itemize}
  \slidesources{\OOPSUnitIISources}
\end{frame}


\begin{frame}{Learning Outcomes}
  \begin{itemize}[<+->]
    \item Write constructors and use \texttt{this(...)} for constructor chaining.
    \item Overload methods correctly and avoid common mistakes.
    \item Use \texttt{static} for class-level state/behavior and explain GC basics.
  \end{itemize}
  \slidesources{\OOPSUnitIISources}
\end{frame}

\begin{frame}{Key Terms to Know}
\small
\begin{columns}[t]
\column{0.48\textwidth}
\begin{itemize}
  \item \textbf{Constructor}: special method to initialize an object (same name as class, no return type).
  \item \textbf{Shadowing}: parameter/local variable has same name as field (e.g., \texttt{name}).
  \item \textbf{\texttt{this}}: reference to the current object.
  \item \textbf{\texttt{this(...)} }: call another constructor in the same class (must be first statement).
\end{itemize}
\column{0.48\textwidth}
\begin{itemize}
  \item \textbf{Overloading}: same method name, different parameter list (compile-time selection).
  \item \textbf{\texttt{static}}: belongs to the class, shared across all objects.
  \item \textbf{Garbage collection}: automatic memory reclamation for unreachable objects.
\end{itemize}
\end{columns}
  \slidesources{\OOPSUnitIISources}
\end{frame}

\begin{frame}{Study Flow for This Lecture}
\begin{enumerate}[<+->]
  \item Refresh the prerequisite bridge before reading new ideas in isolation.
  \item Study the main build in this order: \textit{Constructors and constructor chaining} $\rightarrow$ \textit{\texttt{this} vs \texttt{this(...)}} $\rightarrow$ \textit{Method overloading (rules)} $\rightarrow$ \textit{\texttt{static} keyword}.
  \item Trace the worked example and connect each line back to one concept frame.
  \item Attempt the mini exercises before moving to the recap and exit check.
\end{enumerate}
  \slidesources{\OOPSUnitIISources}
\end{frame}


\section{Key Topics}
\label{sec:concepts}

\begin{frame}{Unit 02: Classes, Inheritance, Packages and Interfaces}
  \begin{itemize}[<+->]
    \item Lecture focus: Constructors, `this`, Overloading, `static`, and GC Overview
  \end{itemize}
  \slidesources{\OOPSUnitIISources}
\end{frame}

\begin{frame}{Today: Roadmap}
  \begin{itemize}[<+->]
    \item Constructors (default/parameterized), constructor chaining, `this()` / `this`
    \item Method overloading (rules + examples)
    \item `static` keyword (fields/methods/blocks) and utility-class pattern
    \item Garbage collection overview (what it is / what it isn't)
  \end{itemize}
  \slidesources{\OOPSUnitIISources}
\end{frame}

\begin{frame}{Constructors (why and how)}
  \begin{itemize}[<+->]
    \item Constructor initializes an object and enforces invariants.
    \item If you define any constructor, Java does \textbf{not} add a default no-arg one automatically.
    \item Use \texttt{this(...)} for constructor chaining (must be first statement).
  \end{itemize}
  \slidesources{\OOPSUnitIISources}
\end{frame}

\begin{frame}{\texttt{this} vs \texttt{this(...)}}
  \begin{itemize}[<+->]
    \item \texttt{this} = current object reference (resolves shadowing: \texttt{this.name = name;}).
    \item \texttt{this(...)} = calls another constructor (reduces duplicate code).
  \end{itemize}
  \slidesources{\OOPSUnitIISources}
\end{frame}

\begin{frame}{Method overloading (rules)}
  \begin{itemize}[<+->]
    \item Same method name, different parameter list (number/type/order).
    \item Not overloading: changing only return type.
    \item Selected at compile time (based on argument types).
  \end{itemize}
  \slidesources{\OOPSUnitIISources}
\end{frame}

\begin{frame}{\texttt{static} keyword}
  \begin{itemize}[<+->]
    \item Static field is shared across all objects (e.g., counters).
    \item Static method belongs to class: \texttt{ClassName.method()}.
    \item Static methods cannot access instance fields directly (no \texttt{this}).
  \end{itemize}
  \slidesources{\OOPSUnitIISources}
\end{frame}

\begin{frame}{Garbage collection (realistic view)}
  \begin{itemize}[<+->]
    \item GC reclaims memory of objects that are no longer reachable.
    \item \texttt{System.gc()} is a suggestion, not a guarantee.
    \item GC does not automatically close files/DB connections (you must close resources).
  \end{itemize}
  \slidesources{\OOPSUnitIISources}
\end{frame}

\begin{frame}{Checkpoint and Reflection}
\begin{block}{Reflection prompt}
Where would \textit{Constructors, `this`, Overloading, `static`, and GC Overview} matter most in Campus Resource Manager, and which design choice would you justify using this lecture's ideas?
\end{block}
\begin{block}{Quick self-check}
Which part needs one more example before you move on: \textit{Constructors and constructor chaining}, \textit{\texttt{this} vs \texttt{this(...)}}, the worked example, or the recap?
\end{block}
  \slidesources{\OOPSUnitIISources}
\end{frame}

\section{Demo / Example}
\label{sec:demo}

\begin{frame}[fragile,shrink=8]{Example: constructors + \texttt{static} counter (1/2)}
\begin{lstlisting}
public class Employee {
    private static int nextId = 1;

    private final int id;
    private final String name;
    private double salary;

    public Employee() { this("NA", 0.0); }

    public Employee(String name, double salary) {
        this.id = nextId++;
        this.name = name;
        this.salary = salary;
    }
\end{lstlisting}
  \slidesources{\OOPSUnitIISources}
\end{frame}

\begin{frame}[fragile]{Example: constructors + \texttt{static} counter (2/2)}
\begin{lstlisting}
// continued...
    public static int getNextId() { return nextId; }
}
\end{lstlisting}
  \slidesources{\OOPSUnitIISources}
\end{frame}

\begin{frame}[fragile]{Example: method overloading}
\begin{lstlisting}
public void raiseSalary(double percent) {
    raiseSalary(percent, 0.0);
}

public void raiseSalary(double percent, double bonus) {
    salary = salary + (salary * percent / 100.0) + bonus;
}
\end{lstlisting}
  \slidesources{\OOPSUnitIISources}
\end{frame}

\begin{frame}{Mini Exercises (2--3 mins)}
  \begin{enumerate}[<+->]
    \item Difference: \texttt{this} vs \texttt{this(...)}?
    \item Can we overload by changing only return type?
    \item Can a static method access instance fields directly?
  \end{enumerate}
  \slidesources{\OOPSUnitIISources}
\end{frame}

\begin{frame}{Watch-outs}
\begin{itemize}[<+->]
  \item Using this and this() as if they solve the same problem.
  \item Trying to overload by return type alone.
  \item Putting per-object data inside a static field.
\end{itemize}
  \slidesources{\OOPSUnitIISources}
\end{frame}

\section{Summary}
\label{sec:summary}

\begin{frame}{Key Takeaways}
  \begin{itemize}[<+->]
    \item Constructors create valid objects and can be chained using \texttt{this(...)}.
    \item Overloading = same name, different parameters (compile-time).
    \item \texttt{static} is class-level; GC is automatic but not controllable exactly.
  \end{itemize}
  \slidesources{\OOPSUnitIISources}
\end{frame}

\begin{frame}{Checkpoint Questions}
  \begin{enumerate}[<+->]
    \item Why is \texttt{this(...)} useful in constructors?
    \item What is one valid and one invalid way to overload a method?
    \item What does garbage collection do (in one sentence)?
  \end{enumerate}
  \slidesources{\OOPSUnitIISources}
\end{frame}

\begin{frame}{Next Study Steps}
\begin{itemize}[<+->]
  \item Revisit the recall references if any older idea still feels weak.
  \item Rework the lecture example without looking at the notes first.
  \item Attempt the self-study practice and then answer the exit questions in full sentences.
  \item Apply the idea once inside Campus Resource Manager to make the concept stick.
\end{itemize}
  \slidesources{\OOPSUnitIISources}
\end{frame}

\begin{frame}{Next Lecture Preview}
  \textbf{Next:} Access Modifiers + Packages + Interface Intro
  \vspace{0.6em}
  \begin{itemize}[<+->]
    \item We will organize code into packages and control visibility.
    \item We will introduce interfaces as ``contracts''.
  \end{itemize}
  \slidesources{\OOPSUnitIISources}
\end{frame}

\end{document}
